\documentclass[11pt,a4paper]{article}

\usepackage[margin=1in]{geometry}
\usepackage[T1]{fontenc}
\usepackage[utf8]{inputenc}
\usepackage{lmodern}
\usepackage{microtype}
\usepackage{amsmath,amssymb,amsfonts,mathtools}
\usepackage{booktabs,longtable,tabularx,array,multirow}
\usepackage{enumitem}
\usepackage{xcolor}
\usepackage[most]{tcolorbox}
\usepackage{hyperref}
\usepackage{cleveref}
\usepackage{setspace}

\hypersetup{
  colorlinks=true,
  linkcolor=blue!55!black,
  citecolor=blue!55!black,
  urlcolor=blue!55!black,
  pdftitle={Exact LNP Transcript and Recursive LaBRADOR Compression for the NIBS Relation}
}

\allowdisplaybreaks
\setlist[itemize]{leftmargin=1.6em,itemsep=0.25em,topsep=0.35em}
\setlist[enumerate]{leftmargin=1.8em,itemsep=0.25em,topsep=0.35em}
\renewcommand{\arraystretch}{1.16}

% -----------------------------------------------------------------------------
% Basic notation
% -----------------------------------------------------------------------------
\newcommand{\Z}{\mathbb{Z}}
\newcommand{\N}{\mathbb{N}}
\newcommand{\F}{\mathbb{F}}
\newcommand{\R}{\mathcal{R}}
\newcommand{\Rp}{R_p}
\newcommand{\Rsig}{R_{q_{\mathrm{sig}}}}
\newcommand{\Romuf}{R_{q_{\mathrm{omuf}}}}
\newcommand{\Rpi}{R_{q_{\pi}}}
\newcommand{\sigmod}{q_{\mathrm{sig}}}
\newcommand{\omufmod}{q_{\mathrm{omuf}}}
\newcommand{\proofmod}{q_{\pi}}
\newcommand{\LNP}{\mathsf{LNP}}
\newcommand{\LNPcom}{\mathsf{LNP}^{\mathsf{com}}}
\newcommand{\LaB}{\mathsf{LaBRADOR}}
\newcommand{\LNPLab}{\mathsf{LNPLab}}
\newcommand{\FS}{\mathsf{FS}}
\newcommand{\Com}{\mathsf{Com}}
\newcommand{\CLaB}{\mathsf{C}_{\mathsf{LaB}}}
\newcommand{\Rej}{\mathsf{Rej}}
\newcommand{\Bin}{\mathsf{Bin}}
\newcommand{\MSIS}{\mathsf{MSIS}}
\newcommand{\MLWE}{\mathsf{MLWE}}
\newcommand{\ct}{\operatorname{ct}}
\newcommand{\op}{\operatorname{op}}
\newcommand{\diag}{\operatorname{diag}}
\newcommand{\Enc}{\operatorname{Enc}}
\newcommand{\Bits}{\operatorname{Bits}}
\newcommand{\Lift}{\operatorname{Lift}}
\newcommand{\powvec}{\operatorname{pow}}
\newcommand{\vecx}[1]{\overrightarrow{#1}}
\newcommand{\norm}[1]{\left\lVert #1\right\rVert}
\newcommand{\abs}[1]{\left\lvert #1\right\rvert}
\newcommand{\inner}[2]{\left\langle #1,#2\right\rangle}
\newcommand{\concat}{\mathbin\Vert}
\newcommand{\getsr}{\xleftarrow{\$}}
\newcommand{\calC}{\mathcal{C}}
\newcommand{\calF}{\mathcal{F}}
\newcommand{\calG}{\mathcal{G}}
\newcommand{\calR}{\mathcal{R}}
\newcommand{\one}{\mathbf{1}}
\newcommand{\zero}{\mathbf{0}}
\newcommand{\transpose}{\mathsf{T}}
\newcommand{\pk}{\mathsf{pk}}

\newtcolorbox{importantbox}[1][]{
  enhanced,breakable,colback=yellow!7,colframe=yellow!45!black,
  fonttitle=\bfseries,title={#1},boxrule=0.7pt,arc=1mm,
  left=1.2mm,right=1.2mm,top=1mm,bottom=1mm
}
\newtcolorbox{definitionbox}[1][]{
  enhanced,breakable,colback=blue!3,colframe=blue!45!black,
  fonttitle=\bfseries,title={#1},boxrule=0.7pt,arc=1mm,
  left=1.2mm,right=1.2mm,top=1mm,bottom=1mm
}
\newtcolorbox{protocolbox}[1][]{
  enhanced,breakable,colback=gray!3,colframe=black!55,
  fonttitle=\bfseries,title={#1},boxrule=0.7pt,arc=1mm,
  left=1.2mm,right=1.2mm,top=1mm,bottom=1mm
}

\title{Exact LNP Transcript and Recursive LaBRADOR Compression\\
for the Lattice-Based NIBS Relation}
\author{Technical appendix for \texttt{Master\_thesis\_Tran}}
\date{August 2026}

\begin{document}
\maketitle
\tableofcontents
\newpage

% =============================================================================
\section{Scope and consistency conventions}
% =============================================================================

This document gives a self-contained specification of the exact lattice proof
layer used for the NIBS relation: the application relation, its embedding into
an LNP exact proof, the complete nine-message LNP transcript, the
composition-friendly compressed transcript, the LaBRADOR boundary relation,
and the recursive LaBRADOR proof.  No construction or security argument that
is unrelated to LNP or LaBRADOR is included.

The presentation follows the LNP framework of Lyubashevsky--Nguyen--Plancon
\cite{LNP22}, the LaBRADOR framework of Beullens--Seiler \cite{BS23}, and the
formal LNP--LaBRADOR composition of del Pino et al. \cite{DPKNRT25}.  The
application equations and numerical parameters are taken from the current
master-thesis draft.

\begin{importantbox}[Consistency decisions used throughout this file]
The current thesis PDF contains several shorthand expressions and dimension
entries that are mutually inconsistent.  A transcript cannot be implemented
until one convention is fixed.  This file uses the following conventions.
\begin{enumerate}
  \item The selection vector has length
  \(\ell=8\cdot 8=64\), because it is indexed as
  \(v=(v_{i,j})_{i,j\in[8]}\).  The value ``28'' in the summary table is not
  used.
  \item The rounding-error vector satisfies \(e_M\in R^{m}\) with \(m=21\).
  Thus its ring-coordinate dimension is \(21\), not \(24\).
  \item The pre-signature witness is the full vector
  \(z=(z_L,z_R)\in R^{2k}\).  Consequently, the preimage equation contains
  both terms \(A_{\mathrm{td}}z_L\) and \(B_vz_R\).  Omitting
  \(A_{\mathrm{td}}z_L\) changes the application relation.
  \item The symbol \(z\) is reserved for the application pre-signature.
  The final LNP responses are written
  \((z_{A,1},z_{A,2},z_u)\), and the LaBRADOR folded opening is written
  \(z_{\mathrm{fold}}\).
  \item The LNP final challenge space and the LaBRADOR folding challenge space
  are different sets.  They are denoted \(\calC_{\LNP}\) and
  \(\calC_{\LaB}\), respectively.
\end{enumerate}
All proof-size figures reproduced from the thesis are therefore labelled
\emph{draft estimates}.  They must be regenerated after the final serialization
of the full vector \(z=(z_L,z_R)\), the quotient/carry block, and all binary
range blocks is fixed.
\end{importantbox}

% =============================================================================
\section{Rings, coefficient embeddings, and automorphisms}
% =============================================================================

Let
\[
  R := \Z[X]/(X^d+1), \qquad d=64,
\]
and, for a positive modulus \(q\), let
\[
  R_q := \Z_q[X]/(X^d+1).
\]
We use three application/proof rings:
\[
  \Rsig=R/\sigmod R,\qquad
  \Romuf=R/\omufmod R,\qquad
  \Rpi=R/\proofmod R.
\]
For \(a\in R_q\), \(\vecx a\in\Z_q^d\) denotes its coefficient vector.  The
same notation is extended componentwise to vectors of ring elements.  The
constant coefficient is denoted by \(\ct(a)\).

The conjugation automorphism used by both LNP and LaBRADOR is
\[
  \sigma:R_q\longrightarrow R_q,\qquad
  \sigma(X)=X^{-1}=-X^{d-1}.
\]
For any \(a,b\in R_q\),
\begin{equation}
  \inner{\vecx a}{\vecx b}
  =\ct\!\left(a\,\sigma(b)\right).
  \label{eq:ct-inner-product}
\end{equation}
For polynomial vectors, the right-hand side is summed over coordinates.  This
identity is the bridge from integer inner products and norm equations to
constant-coefficient quadratic equations in \(R_q\).

For \(x\in\Z_q\), let \(x^*\in\{0,1\}^{\lceil\log_2q\rceil}\) be its binary
expansion, and extend \((\cdot)^*\) coefficientwise to ring elements and ring
vectors.  Let
\[
  \powvec(B):=(1,2,\ldots,2^{\lceil\log_2 B\rceil})
\]
with zero padding as necessary.  Thus a nonnegative integer \(a\le B\) can be
written \(a=\inner{a^*}{\powvec(B)}\).

% =============================================================================
\section{Concrete parameter set}
% =============================================================================

\subsection{Fixed and derived parameters}

\begin{longtable}{@{}p{0.25\textwidth}p{0.18\textwidth}p{0.47\textwidth}@{}}
\toprule
Parameter & Value & Meaning / status \\
\midrule
\endfirsthead
\toprule
Parameter & Value & Meaning / status \\
\midrule
\endhead
Security level & \(\lambda=128\) & Target classical/quantum-safe soundness level. \\
Ring degree & \(d=64\) & \(R=\Z[X]/(X^{64}+1)\). \\
Signature modulus & \(\sigmod=2^{36}-599=68{,}719{,}476{,}137\) & Composite modulus with \(\sigmod=13p\). \\
Rounding modulus & \(p=5{,}286{,}113{,}549\) & Prime factor \(p=\sigmod/13\). \\
OMUF modulus & \(\omufmod=2^{36}-5=68{,}719{,}476{,}731\) & OMUF/partitioning modulus from the application. \\
Proof modulus & \(\proofmod=2^{40}-195=1{,}099{,}511{,}627{,}581\) & Prime, and \(\proofmod\equiv5\pmod 8\). \\
Recipient dimension & \(m=21\) & Number of ring coordinates of \(s_R,e_R,e_M\). \\
Preimage output dimension & \(k=72\) & \(z_L,z_R\in R^k\), hence \(z\in R^{2k}\). \\
Selection length & \(\ell=64\) & \(v=(v_{i,j})_{i,j\in[8]}\). \\
Application bounds & \(B_{\mathrm{pre}}=100\), \(B_s=100\) & Bounds for \(z\) and \(s_R\). \\
Error coefficient bound & \(\norm{e_R}_\infty,\norm{e_M}_\infty\le6\) & Coefficientwise rounding/error bound. \\
Derived error bounds & \(B_r=B_M=220\) & \(\lceil6\sqrt{md}\rceil=220\). \\
JL constants & \(C_1=337\), \(C_2=13/2\) & LNP approximate-range/JL constants. \\
Projection rejection factor & \(\gamma_{\mathrm{proj}}=2.5\) & Gaussian mask multiplier. \\
Projection tail factor & \(t=1.65\) & Chosen above the theoretical value \(1.64\). \\
Final LNP rejection factors & \(\gamma_1=17\), \(\gamma_2=1.2\) & Ajtai-message and randomness response samplers. \\
Constant-term repetitions & \(\tau_{\mathrm{cc}}=4\) & \(\lceil128/\log_2\proofmod\rceil=4\). \\
Commitment rank & \(\kappa_{\mathrm{MSIS}}=2\) & Draft choice for compact Ajtai commitments. \\
Binary/range placeholder & \(\mu_{\mathrm{draft}}=4\) & Value used in the current numerical projection calculation; see the warning below. \\
Projection Gaussian width & \(s_{\mathrm{proj}}=15{,}720\) & Ceiling of the expression in \cref{eq:sproj}. \\
Extracted short-vector bound & \(B_R=81{,}385\) & Ceiling of \(t s_{\mathrm{proj}}\sqrt{d/C_2}\). \\
Verifier projection bound & \(B_{\mathrm{JL}}=207{,}492\) & \(\lceil\sqrt{C_2}B_R\rceil\). \\
Compact commitment size & \(640\) bytes & \(2d\log_2\proofmod/8\). \\
One appended ring element & \(320\) bytes & \(d\log_2\proofmod/8\). \\
Projection response size & \(608\) bytes & \(256\) signed entries encoded in \(19\) bits. \\
\bottomrule
\end{longtable}

The projection standard deviation is computed as
\begin{equation}
  s_{\mathrm{proj}}
  =\left\lceil
  \gamma_{\mathrm{proj}}
  \sqrt{C_1\Bigl(
     B_{\mathrm{pre}}^2+B_s^2+B_r^2+B_M^2
     +(\mu_{\mathrm{proj}}+4)d
  \Bigr)}
  \right\rceil.
  \label{eq:sproj}
\end{equation}
Using the draft value \(\mu_{\mathrm{proj}}=4\) gives
\[
  s_{\mathrm{proj}}=15{,}720,
  \qquad
  B_R=\left\lceil1.65\cdot15{,}720\sqrt{64/(13/2)}\right\rceil
  =81{,}385.
\]

\begin{importantbox}[The value \(\mu_{\mathrm{draft}}=4\) is not a complete serialization count]
The explicit range encoding creates many bit-plane ring elements.  For example,
encoding both \(e_R\) and \(e_M\) coefficientwise with two four-bit
representations creates \(8m=168\) ring elements for each error vector, before
adding quotient/carry bits.  Therefore \(\mu_{\mathrm{proj}}=4\) can only be
used if \(\mu\) denotes four \emph{logical blocks}, or if these bit planes are
not part of the projected short vector.  The implementation must define the
projected serialization explicitly, set \(\mu_{\mathrm{proj}}\) to its actual
ring-coordinate dimension, and recompute \(s_{\mathrm{proj}},B_R\), and all
no-wrap bounds.  The transcript below is written symbolically so that this
recomputation changes parameters, not protocol logic.
\end{importantbox}

\subsection{No-wrap checks at the displayed draft parameters}

With four explicit norm-slack polynomials, the full corrected coordinate count
for the four bounded arithmetic blocks is
\[
  L_{\mathrm{full}}
   =2k+m+m+m+\mu_{\mathrm{draft}}+4
   =144+21+21+21+4+4=215.
\]
The current thesis calculation uses \(L_{\mathrm{draft}}=214\).  Both values
satisfy the no-wrap inequalities.  Using the conservative value
\(L_{\mathrm{full}}=215\), define
\begin{align*}
  T_1&:=B_R^2+\sqrt{L_{\mathrm{full}}d}\,B_R
       <6.634\times10^9,\\
  T_2&:=41dL_{\mathrm{full}}\sqrt{C_2}\,B_R
       <1.171\times10^{11},\\
  T_{\mathrm{pre}}&:=B_R^2+2B_{\mathrm{pre}}^2
       =6{,}623{,}538{,}225,\\
  T_r&:=B_R^2+2B_r^2
       =6{,}623{,}615{,}025,\\
  T_M&:=B_R^2+2B_M^2
       =6{,}623{,}615{,}025,\\
  T_s&:=B_R^2+2B_s^2
       =6{,}623{,}538{,}225.
\end{align*}
Hence
\begin{equation}
  \max\{T_1,T_2,T_{\mathrm{pre}},T_r,T_M,T_s\}
  =T_2<\proofmod,
  \label{eq:no-wrap-summary}
\end{equation}
and \(\proofmod/T_2>9.39\).  The verifier checks
\[
  \norm{z_{\mathrm{JL}}}_2\le
  \sqrt{C_2}B_R<207{,}492.
\]

\subsection{Parameters that are not fixed by the current sources}

The following parameters cannot be read off from the thesis or the cited
papers and must be selected by the implementation/security estimator:
\begin{itemize}
  \item the Ajtai/BDLOP opening dimension \(\kappa_A\) and all opening
  distributions;
  \item the exact coefficient bounds and binary lengths of every quotient/carry
  vector \(\kappa_j\);
  \item the actual projected binary/range dimension
  \(\mu_{\mathrm{proj}}\);
  \item the LNP challenge parameters
  \((\kappa_{\mathrm{ch}},k_{\mathrm{ch}},\eta_{\mathrm{ch}})\);
  \item the per-recursion LaBRADOR commitment ranks, bases
  \((b,b_1,b_2)\), and decomposition lengths \((t_1,t_2)\), unless they are
  generated by the reference parameter-selection code.
\end{itemize}
No concrete security claim should be attached to the numerical proof-size
estimate until these values have been fixed and checked against Module-SIS and
Module-LWE estimators.

% =============================================================================
\section{The application relation to be proved}
% =============================================================================

\subsection{Public instance and witness}

Let the public instance be
\begin{align*}
 x_{\mathrm{app}}:=(&A,A_0,A_1,A_{\mathrm{td}},u,C,B_0,
       (B_{1,i})_{i\in[8]},(B_{2,j})_{j\in[8]},\pk_R,M),
\end{align*}
where all dimensions are those of the NIBS construction.  In particular,
\(A_0,A_1\in\Rsig^{m\times m}\), \(C\) maps the recipient public key to the
preimage-equation output space, and \(A_{\mathrm{td}}\) and \(B_v\) act on the
two halves of the pre-signature vector.

For \(i,j\in[8]\), define the public selected-matrix blocks
\[
  D_{i,j}:=B_{1,i}G^{-1}(B_{2,j}),
\]
and for \(v=(v_{i,j})_{i,j\in[8]}\), define
\begin{equation}
  B_v:=B_0+\sum_{i=1}^{8}\sum_{j=1}^{8}v_{i,j}D_{i,j}.
  \label{eq:Bv}
\end{equation}

The native witness is
\[
  \omega_{\mathrm{nat}}
   =(s_R,e_R,z_L,z_R,v,e_M,s_1,\ldots,s_{\ell}),
\]
with \(s_0:=s_R\),
\[
  s_R,e_R,e_M,s_i\in R^m,\qquad
  z_L,z_R\in R^k,\qquad
  v\in\{0,1\}^{64}\setminus\{0^{64}\}.
\]

\subsection{Native equations over \(R_{q_{\mathrm{sig}}}\)}

The relation accepts if all of the following equations and bounds hold.

\paragraph{Recipient-key equation.}
\begin{equation}
  13\pk_R-A s_R-e_R=0
  \qquad\text{in }\Rsig^m.
  \label{eq:R1-native}
\end{equation}

\paragraph{Selected trapdoor-preimage equation.}
\begin{equation}
  A_{\mathrm{td}}z_L+B_vz_R-u+C(13\pk_R)=0
  \qquad\text{in }\Rsig^k.
  \label{eq:R2-native}
\end{equation}
The implementation convention for \(A_{\mathrm{td}}z_L\) and \(B_vz_R\)
must be the same convention used by \textsf{SampleLeft}.  It yields \(k=72\)
ring-coordinate equations.

\paragraph{PRF recursion.}
For \(i=0,\ldots,\ell-1\),
\begin{equation}
  s_{i+1}-\bigl(A_0+(A_1-A_0)v_i\bigr)s_i=0
  \qquad\text{in }\Rsig^m.
  \label{eq:R3-native}
\end{equation}
These are \(\ell m=64\cdot21=1{,}344\) ring-coordinate equations.

\paragraph{Rounding equation.}
\begin{equation}
  s_{\ell}+e_M-13M=0
  \qquad\text{in }\Rsig^m.
  \label{eq:R4-native}
\end{equation}
This equation is the exact algebraic form of \(M=[s_\ell]_p\), together with
the coefficientwise rounding-error bound on \(e_M\).

\paragraph{Bounds and Boolean conditions.}
\begin{align}
  \norm{z_L\concat z_R}_2&\le B_{\mathrm{pre}}=100,
  &\norm{s_R}_2&\le B_s=100,\label{eq:app-norm1}\\
  \norm{e_R}_\infty&\le6,
  &\norm{e_M}_\infty&\le6,\label{eq:app-norm2}\\
  v_i(v_i-1)&=0\quad(i\in[64]),
  &\sum_{i=1}^{64}v_i&\ge1.\label{eq:v-nonzero}
\end{align}
The coefficientwise error bounds imply
\(\norm{e_R}_2,\norm{e_M}_2\le220\).

\subsection{Equation count before LNP aggregation}

The four equation families contain
\begin{equation}
  N_{\mathrm{full}}
   =m+k+\ell m+m
   =21+72+1{,}344+21
   =1{,}458
  \label{eq:Nfull}
\end{equation}
ring-coordinate equations.  The quotient/carry block introduced below must
therefore contain at least one ring-coordinate quotient for each of these
\(1{,}458\) outputs, unless several outputs are combined by a separately proven
injective encoding.

% =============================================================================
\section{Lifting the application equations to the LNP proof modulus}
% =============================================================================

LNP is executed over \(\Rpi\), whereas
\cref{eq:R1-native,eq:R2-native,eq:R3-native,eq:R4-native} are equations modulo
\(\sigmod\).  It is not sufficient to reinterpret the same expressions modulo
\(\proofmod\).  Every equation must be lifted through a bounded quotient.

For each ring-coordinate equation \(F_j(\omega)=0\pmod{\sigmod}\), choose
centered integer representatives and write
\begin{equation}
  \widetilde F_j(\omega)-\sigmod\,\kappa_j=0
  \qquad\text{over }R,
  \label{eq:integer-lift}
\end{equation}
for a quotient/carry polynomial \(\kappa_j\in R\).  The LNP relation contains
the corresponding proof-ring equation
\begin{equation}
  \widetilde F_j(\omega)-\sigmod\,\kappa_j=0
  \qquad\text{in }\Rpi.
  \label{eq:proof-ring-lift}
\end{equation}
The quotient is represented with a certified signed-binary encoding
\(\rho_{\kappa_j}\), and the relation checks both its reconstruction and the
bitness of that encoding.

\begin{importantbox}[Why quotient bounds are mandatory]
If \(\kappa_j\) is unrestricted, then
\cref{eq:proof-ring-lift} is satisfiable for arbitrary
\(\widetilde F_j(\omega)\) and does not prove the original equation modulo
\(\sigmod\).  The compiler must derive a coefficient bound
\(B_{\kappa,j}\), prove
\(\norm{\kappa_j}_\infty\le B_{\kappa,j}\), and verify that the centered
integer value in \cref{eq:integer-lift} cannot wrap modulo \(\proofmod\).
\end{importantbox}

Let
\[
  \kappa:=(\kappa_1,\ldots,\kappa_{N_{\mathrm{full}}})
\]
be the complete quotient block.  The lifted full-ring equations are denoted
\[
  q_j(s^{(0)})=0\quad\text{in }\Rpi,
  \qquad j\in[N_{\mathrm{full}}],
\]
where \(s^{(0)}\) is the ordered vector of all arithmetic witnesses, binary
encodings, norm slacks, and quotient/carry witnesses used before the
constant-coefficient masks are appended.

% =============================================================================
\section{Commitment layout for the exact LNP proof}
% =============================================================================

\subsection{Ajtai block for small and binary data}

Let
\begin{equation}
  m_{\mathrm{bin}}
   :=(v,\rho_R,\rho_M,\rho_\kappa,\rho_{\mathrm{aux}})
  \label{eq:mbin}
\end{equation}
be the concatenation of:
\begin{itemize}
  \item the selector bits \(v\);
  \item the positive/negative range bits for \(e_R\) and \(e_M\);
  \item all signed quotient/carry bits;
  \item the bits used for the nonzero-selector condition and any other bounded
  auxiliary integer.
\end{itemize}
Sample a short Ajtai randomness \(u_A\) and set
\begin{equation}
  V_A=A_1m_{\mathrm{bin}}+A_2u_A
  \qquad\text{in }\Rpi^{\kappa_{\mathrm{MSIS}}}.
  \label{eq:Ajtai-commitment}
\end{equation}
The matrices \(A_1,A_2\) and the opening dimensions belong to the CRS.

\subsection{One appendable BDLOP randomness}

Sample one short vector \(u_{\mathrm{BD}}\).  The base commitment is
\begin{equation}
  t_A=A_2u_{\mathrm{BD}}.
  \label{eq:BDLOP-base}
\end{equation}
For every appended message \(x\), with an independently sampled public matrix
\(B_x\) of the appropriate dimensions, set
\begin{equation}
  t_x=B_xu_{\mathrm{BD}}+x.
  \label{eq:BDLOP-append}
\end{equation}
All append operations use the \emph{same} \(u_{\mathrm{BD}}\).  The initial
append list is
\begin{equation}
 \mathcal X_0=(s_R,e_R,z_L,z_R,S,e_M,\kappa,
              \beta_{\mathrm{pre}},\beta_s,\beta_r,\beta_M,Y_R),
 \label{eq:initial-appends}
\end{equation}
where \(S=(s_1,\ldots,s_\ell)\), the \(\beta\)'s are exact norm slacks, and
\(Y_R\in\Rpi^4\) packs the \(256\) projection-mask integers.  Later rounds
append \(g_1,\ldots,g_{\tau_{\mathrm{cc}}}\) and the quadratic garbage term
\(g_Q\).

\subsection{Coefficientwise error encoding}

For \(x\in\{e_R,e_M\}\), every coefficient \(x[a][j]\) is certified to lie in
\([-6,6]\) using bits
\(\rho^{x,+}_{a,j,t},\rho^{x,-}_{a,j,t}\in\{0,1\}\), \(t=0,1,2,3\), and
\begin{align}
  x[a][j]+6&=\sum_{t=0}^{3}2^t\rho^{x,+}_{a,j,t},
  \label{eq:range-plus}\\
  6-x[a][j]&=\sum_{t=0}^{3}2^t\rho^{x,-}_{a,j,t}.
  \label{eq:range-minus}
\end{align}
For one \(R^m\) error vector this produces \(8m=168\) bit-plane ring
elements after packing over the degree-\(64\) coefficients.

\subsection{Exact norm inequalities as equalities}

For each bounded block
\[
  x\in\{z_L\concat z_R,s_R,e_R,e_M\}
\]
with bound \(B_x\), introduce a binary slack \(\beta_x\) representing
\(B_x^2-\norm{x}_2^2\).  The relation checks
\begin{align}
  \norm{x}_2^2+\inner{\vecx{\beta_x}}{\powvec(B_x^2)}&=B_x^2,
  \label{eq:norm-slack}\\
  \inner{\vecx{\beta_x}}{\vecx{\beta_x}-\one}&=0.
  \label{eq:norm-slack-bits}
\end{align}
By \cref{eq:ct-inner-product}, these become constant-coefficient quadratic
relations in \(\Rpi\).  The no-wrap proof ensures that the equalities hold over
\(\Z\), not merely modulo \(\proofmod\).

The bitness of every bit vector \(b\), including \(v\), range bits, quotient
bits, norm slacks, and nonzero-selector auxiliary bits, is expressed as
\begin{equation}
  \inner{\vecx b}{\vecx b-\one}=0.
  \label{eq:generic-bitness}
\end{equation}
To enforce \(v\ne0\), introduce six bits \(\nu_0,\ldots,\nu_5\) and check
\begin{equation}
  \sum_{i=1}^{64}v_i
  =1+\sum_{t=0}^{5}2^t\nu_t.
  \label{eq:nonzero-selector}
\end{equation}

% =============================================================================
\section{LNP challenge spaces}
% =============================================================================

The four verifier messages in the exact LNP transcript have distinct roles and
must use domain-separated random-oracle expansions after Fiat--Shamir.

\subsection{JL projection challenge}

Let \(D_{\mathrm{proj}}\) be the integer coefficient dimension of the exact
projected serialization \(x_{\mathrm{short}}\).  Each entry of
\(R_{\mathrm{JL}}\) is sampled from the centered binomial distribution
\[
  \Bin_1:\qquad
  \Pr[-1]=\Pr[1]=\tfrac14,\qquad \Pr[0]=\tfrac12.
\]
The verifier samples
\begin{equation}
  R_{\mathrm{JL}}\getsr
  \Bin_1^{256\times D_{\mathrm{proj}}}.
  \label{eq:JL-challenge}
\end{equation}
Using uniform \(\{0,1\}\) entries instead would change the mean and requires a
new proof of the displayed JL constants.

\subsection{Constant-coefficient aggregation challenge}

For \(N_{\mathrm{cc}}\) constant-coefficient equations, the verifier samples
\begin{equation}
  \Gamma=(\gamma_{a,j})
  \getsr\Z_{\proofmod}^{\tau_{\mathrm{cc}}\times N_{\mathrm{cc}}},
  \qquad \tau_{\mathrm{cc}}=4.
  \label{eq:Gamma-challenge}
\end{equation}

\subsection{Full equation aggregation challenge}

After the masked constant-coefficient equations are fixed, the verifier samples
\begin{equation}
  \boldsymbol\mu\getsr
  \Rpi^{N_{\mathrm{full}}+\tau_{\mathrm{cc}}}.
  \label{eq:mu-challenge}
\end{equation}
This challenge is sampled before the prover constructs the final quadratic
proof masks.

\subsection{Final LNP challenge space}

Let
\[
  S_\kappa:=\{c\in R:\norm{c}_\infty\le\kappa\},
  \qquad
  S_\kappa^\sigma:=\{c\in S_\kappa:\sigma(c)=c\}.
\]
For a power of two \(k_{\mathrm{ch}}\), define the LNP challenge set
\begin{equation}
 \calC_{\LNP}:=
 \left\{
 c\in S_{\kappa_{\mathrm{ch}}}^{\sigma}:
 \sqrt[2k_{\mathrm{ch}}]{
   \norm{\sigma(c^{k_{\mathrm{ch}}})c^{k_{\mathrm{ch}}}}_1
 }
 \le\eta_{\mathrm{ch}}
 \right\}.
 \label{eq:LNP-challenge-space}
\end{equation}
Because \(\sigma(c)=c\), a challenge has the form
\begin{equation}
 c=c_0+\sum_{i=1}^{d/2-1}c_i(X^i-X^{d-i}),
 \label{eq:sigma-fixed-c}
\end{equation}
with \(d/2=32\) independent coefficients.  The bound in
\cref{eq:LNP-challenge-space} controls the operator norm of multiplication by
\(c\), and therefore the Gaussian widths used in rejection sampling.

Since \(\proofmod\equiv5\pmod8\), a nonzero polynomial fixed by \(\sigma\) is
invertible under the short-coefficient condition used by LNP.  In extraction,
we need every nonzero difference \(c-c'\) of distinct accepted challenges to
be invertible.

\begin{importantbox}[A concrete \(d=64\) LNP challenge is still an implementation parameter]
The example \((d,\kappa,\eta)=(128,2,59)\) in \cite{LNP22} cannot simply be
copied to \(d=64\), because \(S_2^\sigma\) then contains only about \(2^{74}\)
polynomials.  A natural starting point is
\(\kappa_{\mathrm{ch}}=8\), for which
\((2\kappa_{\mathrm{ch}}+1)^{32}=17^{32}>2^{130.7}\).
The value \(\eta_{\mathrm{ch}}\) must be selected as a high empirical quantile
of the exact bound in \cref{eq:LNP-challenge-space}, and all response widths
and MSIS bounds must then be recomputed.  This calibration is not present in
the current thesis and is therefore not silently assigned a final value here.
\end{importantbox}

% =============================================================================
\section{Projection layer and constant-coefficient embedding}
% =============================================================================

\subsection{Projected short vector}

Let \(x_{\mathrm{short}}\) be the exact ordered concatenation of every witness
block whose integer smallness is used to rule out wrap-around.  At minimum it
contains
\[
  z_L,z_R,s_R,e_R,e_M,
  \beta_{\mathrm{pre}},\beta_s,\beta_r,\beta_M,
\]
and it must additionally contain any binary/range or quotient blocks whose
smallness is not already guaranteed by a binding relation.  Let
\(D_{\mathrm{proj}}\) be its coefficient dimension.

The prover samples
\[
  y_R\getsr D_{\Z^{256},s_{\mathrm{proj}}}
\]
and packs the \(256\) integers into four elements
\(Y_R\in\Rpi^4\), already appended as
\(t_{Y_R}=B_{Y_R}u_{\mathrm{BD}}+Y_R\).
After receiving \(R_{\mathrm{JL}}\), it computes
\begin{equation}
  z_{\mathrm{JL}}
  =R_{\mathrm{JL}}\vecx{x_{\mathrm{short}}}+y_R
  \quad\text{in }\Z_{\proofmod}^{256}.
  \label{eq:projection-response}
\end{equation}
It applies the selected LNP rejection sampler and aborts on rejection.  The
verifier checks
\begin{equation}
  \norm{z_{\mathrm{JL}}}_2\le B_{\mathrm{JL}}.
  \label{eq:projection-norm-check}
\end{equation}

For each row \(r_a\) of \(R_{\mathrm{JL}}\), the relation includes a
constant-coefficient equation proving
\[
  z_{\mathrm{JL},a}
  =\inner{r_a}{\vecx{x_{\mathrm{short}}}}+y_{R,a}.
\]
Thus the projection contributes \(256\) constant-coefficient equations, but
these are subsequently aggregated and do not cause a linear increase in the
final LNP response size.

\subsection{Constant-coefficient masks}

Let
\[
  F_1(s^{(0)}),\ldots,F_{N_{\mathrm{cc}}}(s^{(0)})
\]
be all equations for which only the constant coefficient must vanish.  This
family contains the exact-norm equations, all bitness and range equations,
certified quotient-bound equations, the nonzero-selector equation, and the
\(256\) projection-consistency equations.

For \(a\in[\tau_{\mathrm{cc}}]\), the prover samples
\[
  g_a\getsr\{g\in\Rpi:\ct(g)=0\}
\]
and appends
\[
  t_{g_a}=B_{g_a}u_{\mathrm{BD}}+g_a.
\]
After receiving \(\Gamma\), it computes
\begin{equation}
  h_a=g_a+\sum_{j=1}^{N_{\mathrm{cc}}}
       \gamma_{a,j}F_j(s^{(0)}).
  \label{eq:ha}
\end{equation}
The verifier later checks \(\ct(h_a)=0\) for every \(a\in[4]\).

% =============================================================================
\section{Aggregation into one quadratic equation}
% =============================================================================

Let
\[
  q_1(s^{(0)}),\ldots,q_{N_{\mathrm{full}}}(s^{(0)})
\]
be the lifted full-ring equations and let
\(\bar q_a(s^{(0)},g_a,h_a)=0\) be the full-ring equation corresponding to
\cref{eq:ha}.  After receiving \(\boldsymbol\mu\), define
\begin{equation}
 Q_{\boldsymbol\mu}(s_Q)
  :=\sum_{j=1}^{N_{\mathrm{full}}}\mu_jq_j(s^{(0)})
   +\sum_{a=1}^{\tau_{\mathrm{cc}}}
       \mu_{N_{\mathrm{full}}+a}\bar q_a(s^{(0)},g_a,h_a),
 \label{eq:Qmu}
\end{equation}
where
\[
  s_Q:=(s^{(0)},g_1,\ldots,g_{\tau_{\mathrm{cc}}}),
  \qquad
  \bar s_Q:=(s_Q,\sigma(s_Q)).
\]
Every term in \cref{eq:Qmu} has degree at most two.  Therefore the prover and
verifier can derive public/transcript-dependent \(E,e,e_0\) such that
\begin{equation}
  Q_{\boldsymbol\mu}(s_Q)
  =\bar s_Q^{\transpose}E\bar s_Q+e^{\transpose}\bar s_Q+e_0.
  \label{eq:quadratic-matrix-form}
\end{equation}
The desired aggregated relation is
\begin{equation}
  \bar s_Q^{\transpose}E\bar s_Q+e^{\transpose}\bar s_Q+e_0=0
  \qquad\text{in }\Rpi.
  \label{eq:one-quadratic-relation}
\end{equation}

\subsection{Factorization only after the aggregation challenge}

The quadratic circuit should be generated from the already aggregated
\(Q_{\boldsymbol\mu}\), not by first compiling every source equation and then
summing the resulting products.

For the selected-matrix block, the \(72\) output equations aggregate as
\[
 \sum_{a=1}^{72}\mu_aQ^{(2)}_a
 =L^{(2)}_{\boldsymbol\mu}(z)
  +\sum_{r=1}^{64}v_r
    \left(\sum_{a=1}^{72}\mu_aD_{r,a}(z_R)\right),
\]
so this block requires one ring multiplication per selector: \(64\), not
\(64\cdot72\).

For PRF round \(i\),
\[
 \sum_{a=1}^{21}\mu_{i,a}Q^{(3)}_{i,a}
 =L^{(3)}_{\boldsymbol\mu,i}(s_i,s_{i+1},\kappa)
  -v_i\left(\sum_{a=1}^{21}\mu_{i,a}M_{i,a}(s_i)\right),
\]
so each PRF round requires one ring multiplication.  Across all \(64\) rounds,
\(64\cdot21=1{,}344\) source products collapse to \(64\) aggregated products.
This is an exact algebraic refactoring of the same polynomial
\(Q_{\boldsymbol\mu}\); it does not modify an LNP challenge or soundness
parameter.

% =============================================================================
\section{Final quadratic proof and complete nine-message LNP transcript}
% =============================================================================

\subsection{Masks for the two commitment systems}

The prover samples
\begin{align*}
 y_{A,1}&\getsr D_{s_1}^{\dim(m_{\mathrm{bin}})},
 &y_{A,2}&\getsr D_{s_2}^{\kappa_A},
 &y_u&\getsr D_{s_2}^{\kappa_A}.
\end{align*}
It computes
\begin{align}
  w_A&:=A_1y_{A,1}+A_2y_{A,2},
  \label{eq:wA}\\
  w_u&:=A_2y_u.
  \label{eq:wu}
\end{align}

Construct a mask vector \(y_Q\) having the same coordinate layout as \(s_Q\):
\begin{itemize}
  \item the Ajtai-message coordinate uses \(y_{A,1}\);
  \item an appended BDLOP message \(x\), committed as
  \(t_x=B_xu_{\mathrm{BD}}+x\), uses the mask \(-B_xy_u\).
\end{itemize}
Set \(\bar y_Q=(y_Q,\sigma(y_Q))\), and define
\begin{align}
  g_Q&:=\bar s_Q^{\transpose}(E+E^{\transpose})\bar y_Q
       +e^{\transpose}\bar y_Q,
  \label{eq:gQ}\\
  t_Q&:=B_Qu_{\mathrm{BD}}+g_Q,
  \label{eq:tQ}\\
  v_Q&:=\bar y_Q^{\transpose}E\bar y_Q+B_Qy_u.
  \label{eq:vQ}
\end{align}

\subsection{Final challenge and responses}

After receiving \(c\getsr\calC_{\LNP}\), the prover computes
\begin{align}
  z_{A,1}&:=c\,m_{\mathrm{bin}}+y_{A,1},
  \label{eq:zA1}\\
  z_{A,2}&:=c\,u_A+y_{A,2},
  \label{eq:zA2}\\
  z_u&:=c\,u_{\mathrm{BD}}+y_u.
  \label{eq:zu}
\end{align}
It runs the designated rejection samplers with factors \(\gamma_1=17\) and
\(\gamma_2=1.2\), and aborts if either sampler rejects.

For each BDLOP-appended message \(x\), the verifier derives the corresponding
masked coordinate
\begin{equation}
  z_x:=ct_x-B_xz_u=cx-B_xy_u=cx+y_x.
  \label{eq:derive-zx}
\end{equation}
Combining these coordinates with \(z_{A,1}\) produces
\(z_Q=cs_Q+y_Q\), and
\(\bar z_Q=(z_Q,\sigma(z_Q))\).

\subsection{Nine-message interactive transcript}

\begin{protocolbox}[Exact interactive LNP transcript for the NIBS relation]
\begin{longtable}{@{}p{0.07\textwidth}p{0.09\textwidth}p{0.76\textwidth}@{}}
\toprule
Move & Sender & Message and computation \\
\midrule
1 & Prover &
Send
\[
 a_1=(V_A,t_A,t_{s_R},t_{e_R},t_{z_L},t_{z_R},t_S,t_{e_M},
      t_\kappa,t_{\beta_{\mathrm{pre}}},t_{\beta_s},t_{\beta_r},t_{\beta_M},t_{Y_R}).
\]
All BDLOP values use the same opening randomness \(u_{\mathrm{BD}}\). \\
\addlinespace
2 & Verifier &
Sample and send
\(R_{\mathrm{JL}}\getsr\Bin_1^{256\times D_{\mathrm{proj}}}\). \\
\addlinespace
3 & Prover &
Compute \(z_{\mathrm{JL}}\) as in \cref{eq:projection-response}, apply rejection
sampling, append \(g_1,\ldots,g_4\), and send
\[
 a_2=(z_{\mathrm{JL}},t_{g_1},\ldots,t_{g_4}).
\] \\
\addlinespace
4 & Verifier &
Sample and send
\(\Gamma\getsr\Z_{\proofmod}^{4\times N_{\mathrm{cc}}}\). \\
\addlinespace
5 & Prover &
Compute \(h_a\) by \cref{eq:ha} and send
\[
 a_3=(h_1,h_2,h_3,h_4).
\] \\
\addlinespace
6 & Verifier &
Sample and send
\(\boldsymbol\mu\getsr\Rpi^{N_{\mathrm{full}}+4}\). \\
\addlinespace
7 & Prover &
Derive \(E,e,e_0\), sample the quadratic-proof masks, compute
\(w_A,w_u,t_Q,v_Q\), and send
\[
 a_4=(w_A,w_u,t_Q,v_Q).
\] \\
\addlinespace
8 & Verifier &
Sample and send \(c\getsr\calC_{\LNP}\). \\
\addlinespace
9 & Prover &
Compute \((z_{A,1},z_{A,2},z_u)\), apply rejection sampling, and send
\[
 a_5=(z_{A,1},z_{A,2},z_u).
\] \\
\bottomrule
\end{longtable}
\end{protocolbox}

\subsection{Complete verifier checks}

The verifier accepts if all of the following checks pass.

\paragraph{Projection and constant coefficients.}
\begin{align}
  \norm{z_{\mathrm{JL}}}_2&\le B_{\mathrm{JL}},
  \label{eq:V-proj}\\
  \ct(h_a)&=0\qquad(a\in[4]).
  \label{eq:V-ct}
\end{align}

\paragraph{Ajtai opening equation.}
\begin{equation}
  f_1:=cV_A+w_A-A_1z_{A,1}-A_2z_{A,2}=0.
  \label{eq:f1}
\end{equation}

\paragraph{BDLOP base-opening equation.}
\begin{equation}
  f_2:=ct_A+w_u-A_2z_u=0.
  \label{eq:f2}
\end{equation}

\paragraph{Quadratic equation.}
Define
\begin{equation}
  f:=ct_Q-B_Qz_u=cg_Q-B_Qy_u.
  \label{eq:f-correction}
\end{equation}
Then check
\begin{equation}
  f_3:=
  \bar z_Q^{\transpose}E\bar z_Q
  +c\,e^{\transpose}\bar z_Q
  +c^2e_0-f-v_Q=0.
  \label{eq:f3}
\end{equation}
Indeed, using \(z_Q=cs_Q+y_Q\) and
\cref{eq:one-quadratic-relation}, the left-hand side expands to zero.

\paragraph{Response norm.}
Check
\begin{equation}
  \norm{z_{A,1}\concat z_{A,2}\concat z_u}_2
  \le B_{\mathrm{resp}},
  \label{eq:response-bound}
\end{equation}
where \(B_{\mathrm{resp}}\) is derived from the final selected challenge norm,
Gaussian widths, response dimensions, and the required tail probability.  It is
not the same quantity as the application bound \(B_{\mathrm{pre}}=100\).

% =============================================================================
\section{Fiat--Shamir form and transcript serialization}
% =============================================================================

Use independent domain-separated expanders
\(H_{\mathrm{JL}},H_{\mathrm{cc}},H_{\mathrm{agg}},H_{\LNP}\).  For a context
string
\[
  \mathsf{ctx}=\mathsf{encode}(\mathsf{protocol\_id},\mathsf{crs\_id},x_{\mathrm{app}}),
\]
define
\begin{align*}
 R_{\mathrm{JL}}
   &=\mathsf{Expand}_{\mathrm{JL}}
     \bigl(H_{\mathrm{JL}}(\mathsf{ctx},a_1)\bigr),\\
 \Gamma
   &=\mathsf{Expand}_{\mathrm{cc}}
     \bigl(H_{\mathrm{cc}}(\mathsf{ctx},a_1,R_{\mathrm{JL}},a_2)\bigr),\\
 \boldsymbol\mu
   &=\mathsf{Expand}_{\mathrm{agg}}
     \bigl(H_{\mathrm{agg}}(\mathsf{ctx},a_1,R_{\mathrm{JL}},a_2,
       \Gamma,a_3)\bigr),\\
 c
   &=\mathsf{Sample}_{\calC_{\LNP}}
     \bigl(H_{\LNP}(\mathsf{ctx},a_1,R_{\mathrm{JL}},a_2,
       \Gamma,a_3,\boldsymbol\mu,a_4)\bigr).
\end{align*}
All vectors, matrices, signs, dimensions, and ring coefficients must have a
canonical byte encoding.  Challenges are recomputed by the verifier and are
not trusted as independent proof fields.

% =============================================================================
\section{Composition-friendly LNP with compressed intermediate flows}
% =============================================================================

The uncompressed messages \(a_1,a_3,a_4\) scale with the statement and witness.
Following \cite{DPKNRT25}, replace each by a compact, non-hiding Ajtai
commitment to its binary decomposition.

\subsection{Compressing commitment}

For a ring-vector message \(m\in\Rpi^n\), define
\begin{equation}
  \CLaB.\Com(m):=A_{\CLaB}\,m^*,
  \qquad
  A_{\CLaB}\in\Rpi^{\kappa_{\mathrm{MSIS}}\times n\lceil\log_2\proofmod\rceil}.
  \label{eq:CLaB}
\end{equation}
This commitment is deterministic and non-hiding.  It is used only to bind a
large transcript flow to the value opened inside LaBRADOR.  For
\(\kappa_{\mathrm{MSIS}}=2,d=64,\log_2\proofmod=40\), its output has size
\(640\) bytes, independent of \(n\).

To avoid the naming collision between the Ajtai message in
\cref{eq:mbin} and the first compressed transcript flow, define
\begin{align*}
 m_{\mathrm{flow},A}&:=(t_A,t_{s_R},t_{e_R},t_{z_L},t_{z_R},t_S,t_{e_M},
  t_\kappa,t_{\beta_{\mathrm{pre}}},t_{\beta_s},t_{\beta_r},t_{\beta_M},t_{Y_R}),\\
 m_{\mathrm{flow},h}&:=(h_1,h_2,h_3,h_4),\\
 m_{\mathrm{flow},w}&:=(w_A,w_u,t_Q,v_Q),
\end{align*}
and
\begin{equation}
  x_A:=\CLaB.\Com(m_{\mathrm{flow},A}),\quad
  x_h:=\CLaB.\Com(m_{\mathrm{flow},h}),\quad
  x_w:=\CLaB.\Com(m_{\mathrm{flow},w}).
  \label{eq:three-flow-commitments}
\end{equation}

The public LNP transcript now sends \((V_A,x_A)\),
\((z_{\mathrm{JL}},t_{g_1},\ldots,t_{g_4})\), \(x_h\), and \(x_w\), while
the openings of \(x_A,x_h,x_w\) are moved into the final witness.

\subsection{Complete LaBRADOR boundary relation}

Let \(x_{\mathrm{bdry}}\) be the public view up to the final LNP response,
including
\[
  x_{\mathrm{app}},V_A,x_A,z_{\mathrm{JL}},(t_{g_a})_{a\in[4]},x_h,x_w
\]
and all Fiat--Shamir challenges recomputed from these fields.  Let
\begin{align*}
 w_{\mathrm{bdry}}:=(&m_{\mathrm{flow},A},m_{\mathrm{flow},h},m_{\mathrm{flow},w},
 z_{A,1},z_{A,2},z_u,\\
 &m_{\mathrm{flow},A}^*,m_{\mathrm{flow},h}^*,m_{\mathrm{flow},w}^*).
\end{align*}
The complete relation \(\calR_{\mathrm{bdry}}\) accepts if and only if:
\begin{enumerate}
  \item all three equations in \cref{eq:three-flow-commitments} hold and all
  supplied decompositions are binary and reconstruct the opened messages;
  \item the opened messages parse with the exact dimensions fixed by the CRS;
  \item the projection norm check \cref{eq:V-proj} and all four constant-term
  checks \cref{eq:V-ct} hold;
  \item \(E,e,e_0,z_Q\) are recomputed from the opened transcript exactly as in
  the LNP verifier;
  \item \(f_1=f_2=f_3=0\) as in
  \cref{eq:f1,eq:f2,eq:f3};
  \item the response norm check \cref{eq:response-bound} holds.
\end{enumerate}
In set notation,
\begin{equation}
 \calR_{\mathrm{bdry}}
 =\left\{(x_{\mathrm{bdry}},w_{\mathrm{bdry}}):
 \begin{array}{l}
 \mathsf{Open}_{\CLaB}(x_A,m_{\mathrm{flow},A})=1,\\
 \mathsf{Open}_{\CLaB}(x_h,m_{\mathrm{flow},h})=1,\\
 \mathsf{Open}_{\CLaB}(x_w,m_{\mathrm{flow},w})=1,\\
 \mathsf{LNPVerifyAll}(x_{\mathrm{bdry}},w_{\mathrm{bdry}})=1
 \end{array}
 \right\}.
 \label{eq:boundary-relation}
\end{equation}
The abbreviated relation containing only \(f_1,f_2,f_3\), the response norm,
and the three compact openings is sufficient only if the projection,
constant-term, parsing, and challenge-recomputation checks are enforced
elsewhere.  The complete relation above leaves no such implicit checks.

% =============================================================================
\section{LaBRADOR principal relation}
% =============================================================================

Let \(n\ge1\), \(r\ge1\), and \(\beta>0\).  A LaBRADOR statement consists of
two families of quadratic dot-product functions:
\begin{align*}
 f(s_1,\ldots,s_r)
 &=\sum_{i,j=1}^{r}a_{i,j}\inner{s_i}{s_j}
   +\sum_{i=1}^{r}\inner{\varphi_i}{s_i}-b,
   &&f\in\calF,\\
 f^0(s_1,\ldots,s_r)
 &=\sum_{i,j=1}^{r}a^0_{i,j}\inner{s_i}{s_j}
   +\sum_{i=1}^{r}\inner{\varphi_i^0}{s_i}-b_0,
   &&f^0\in\calF^0.
\end{align*}
The principal relation is
\begin{equation}
 \calR_{\LaB}:=
 \left\{
 ((\calF,\calF^0,\beta),(s_1,\ldots,s_r)):
 \begin{array}{l}
 f(s_1,\ldots,s_r)=0\quad\forall f\in\calF,\\
 \ct(f^0(s_1,\ldots,s_r))=0\quad\forall f^0\in\calF^0,\\
 \sum_{i=1}^{r}\norm{s_i}_2^2\le\beta^2
 \end{array}
 \right\}.
 \label{eq:LaB-principal}
\end{equation}

The binary R1CS obtained from \cref{eq:boundary-relation} is reduced to an
instance of \cref{eq:LaB-principal}.  Binary coefficient constraints,
compact-commitment equations, linear verifier equations, and the final
quadratic equation are all expressible as full-ring or constant-coefficient
dot-product constraints.  The LaBRADOR proof then recursively proves knowledge
of a short solution to this principal relation.

% =============================================================================
\section{LaBRADOR challenge space}
% =============================================================================

LaBRADOR requires a challenge set \(\calC_{\LaB}\subset\Rpi\) such that every
difference of two distinct challenges is invertible and
\begin{equation}
  \norm{c}_2^2\le\tau_{\LaB},
  \qquad
  \norm{c}_{\op}\le T_{\op}
  \qquad(c\in\calC_{\LaB}).
  \label{eq:LaB-challenge-bounds}
\end{equation}
For the degree-\(64\) instantiation of \cite{BS23}, a challenge has
\begin{itemize}
  \item \(23\) zero coefficients;
  \item \(31\) coefficients in \(\{\pm1\}\);
  \item \(10\) coefficients in \(\{\pm2\}\).
\end{itemize}
Thus
\[
  \tau_{\LaB}:=\norm{c}_2^2=31+4\cdot10=71,
  \qquad \norm{c}_2=\sqrt{71}.
\]
Rejection sampling retains only challenges with \(\norm{c}_{\op}<15\), so use
\(T_{\op}=15\).  The set contains more than \(2^{128}\) elements, and the
short-difference condition gives invertibility in the selected ring.

% =============================================================================
\section{One nine-message LaBRADOR recursion level}
% =============================================================================

Suppose the current principal-relation witness is
\(s_1,\ldots,s_r\in\Rpi^n\).  Let
\(A\in\Rpi^{\kappa\times n}\) be the inner commitment matrix.  Let
\((b_1,t_1)\) and \((b_2,t_2)\) be the decomposition parameters for arbitrary
and short garbage values, respectively.

\subsection{Preparatory values}

For every \(i\in[r]\), compute the inner commitment
\[
  t_i=As_i
    =\sum_{a=0}^{t_1-1}t_i^{(a)}b_1^a.
\]
For every \(1\le i\le j\le r\), compute
\begin{align*}
 g_{i,j}&:=\inner{s_i}{s_j}
        =\sum_{a=0}^{t_2-1}g_{i,j}^{(a)}b_2^a,\\
 h_{i,j}&:=\frac12\left(
          \inner{\varphi_i}{s_j}+\inner{\varphi_j}{s_i}
        \right)
        =\sum_{a=0}^{t_1-1}h_{i,j}^{(a)}b_1^a.
\end{align*}
The first outer commitment is
\begin{equation}
 u_1=
 \sum_{i=1}^{r}\sum_{a=0}^{t_1-1}B_{i,a}t_i^{(a)}
 +\sum_{1\le i\le j\le r}\sum_{a=0}^{t_2-1}C_{i,j,a}g_{i,j}^{(a)}.
 \label{eq:LaB-u1}
\end{equation}
The second outer commitment, formed later, is
\begin{equation}
 u_2=
 \sum_{1\le i\le j\le r}\sum_{a=0}^{t_1-1}D_{i,j,a}h_{i,j}^{(a)}.
 \label{eq:LaB-u2}
\end{equation}

\subsection{Nine messages}

\begin{protocolbox}[One LaBRADOR base-protocol execution]
\begin{longtable}{@{}p{0.07\textwidth}p{0.09\textwidth}p{0.76\textwidth}@{}}
\toprule
Move & Sender & Message and computation \\
\midrule
1 & Prover & Send the first outer commitment \(u_1\) from
\cref{eq:LaB-u1}. \\
\addlinespace
2 & Verifier & For every \(i\in[r]\), sample a JL matrix
\(\Pi_i\getsr\Bin_1^{256\times nd}\) and send
\((\Pi_1,\ldots,\Pi_r)\). \\
\addlinespace
3 & Prover & Compute
\[
 p=\sum_{i=1}^{r}\Pi_i\vecx{s_i}\in\Z_{\proofmod}^{256}
\]
and send \(p\).  The verifier immediately checks the prescribed projection
norm. \\
\addlinespace
4 & Verifier & Let \(L=\abs{\calF^0}\) and
\(q_0=\lceil128/\log_2\proofmod\rceil\).  Sample and send
\[
 (\psi^{(a)},\omega^{(a)})
 \getsr\Z_{\proofmod}^{L}\times\Z_{\proofmod}^{256},
 \qquad a\in[q_0].
\] \\
\addlinespace
5 & Prover & Aggregate the \(\calF^0\) relations and projection equations into
full-ring equations \(f^{\prime\prime(a)}\).  Extend the required constant
terms to full polynomials \(b^{\prime\prime(a)}\), and send
\((b^{\prime\prime(1)},\ldots,b^{\prime\prime(q_0)})\). \\
\addlinespace
6 & Verifier & Sample and send
\[
 \alpha\getsr\Rpi^{\abs{\calF}},
 \qquad
 \beta\getsr\Rpi^{q_0}.
\]
These challenges aggregate all full-ring functions into one function
\(F\). \\
\addlinespace
7 & Prover & Compute the linear garbage values \(h_{i,j}\), their
base-\(b_1\) decompositions, form \(u_2\) as in \cref{eq:LaB-u2}, and send
\(u_2\). \\
\addlinespace
8 & Verifier & Sample and send independently
\(c_1,\ldots,c_r\getsr\calC_{\LaB}\). \\
\addlinespace
9 & Prover & Send
\[
 z_{\mathrm{fold}}=\sum_{i=1}^{r}c_is_i
\]
together with all decomposed inner commitments and garbage openings
\((t_i^{(a)}),(g_{i,j}^{(a)}),(h_{i,j}^{(a)})\). \\
\bottomrule
\end{longtable}
\end{protocolbox}

\subsection{Verification equations}

Let
\[
 F(s_1,\ldots,s_r)
 =\sum_{i,j=1}^{r}a_{i,j}\inner{s_i}{s_j}
  +\sum_{i=1}^{r}\inner{\varphi_i}{s_i}-b
\]
be the final aggregate.  The verifier checks:
\begin{enumerate}
  \item all constant coefficients of \(b^{\prime\prime(a)}\) are the values
  determined by \(p,\psi^{(a)},\omega^{(a)}\);
  \item symmetry \(g_{i,j}=g_{j,i}\) and \(h_{i,j}=h_{j,i}\);
  \item every base decomposition reconstructs its parent value and all low
  digits are in their centered intervals;
  \item the consolidated target norm bound
  \begin{equation}
    \norm{z_{\mathrm{fold}}^{(0)}}_2^2
    +\norm{z_{\mathrm{fold}}^{(1)}}_2^2
    +\norm{t}_2^2+\norm{g}_2^2+\norm{h}_2^2
    \le(\beta')^2;
    \label{eq:LaB-consolidated-norm}
  \end{equation}
  \item the amortized opening equation
  \begin{equation}
    A z_{\mathrm{fold}}=\sum_{i=1}^{r}c_it_i;
    \label{eq:LaB-amortized-opening}
  \end{equation}
  \item the three garbage/dot-product identities
  \begin{align}
   \inner{z_{\mathrm{fold}}}{z_{\mathrm{fold}}}
     &=\sum_{i,j=1}^{r}g_{i,j}c_ic_j,
     \label{eq:LaB-g1}\\
   \sum_{i=1}^{r}\inner{\varphi_i}{z_{\mathrm{fold}}}c_i
     &=\sum_{i,j=1}^{r}h_{i,j}c_ic_j,
     \label{eq:LaB-g2}\\
   \sum_{i,j=1}^{r}a_{i,j}g_{i,j}+\sum_{i=1}^{r}h_{i,i}-b
     &=0;
     \label{eq:LaB-g3}
  \end{align}
  \item both outer commitment opening equations
  \cref{eq:LaB-u1,eq:LaB-u2}.
\end{enumerate}
These checks are again dot-product constraints plus one norm relation, so the
last prover message is a witness for another instance of
\cref{eq:LaB-principal}.

% =============================================================================
\section{Recursion, decomposition, and norm propagation}
% =============================================================================

Write
\begin{equation}
  z_{\mathrm{fold}}
  =z_{\mathrm{fold}}^{(0)}+b z_{\mathrm{fold}}^{(1)}
  \label{eq:z-decomposition}
\end{equation}
using centered residues modulo a base \(b\ge2\).  Concatenate
\[
  v_{\mathrm{aux}}:=t\concat g\concat h\in\Rpi^{m_{\mathrm{aux}}},
  \qquad
  m_{\mathrm{aux}}
  =rt_1\kappa+(t_1+t_2)\frac{r^2+r}{2}.
\]
Split each of \(z_{\mathrm{fold}}^{(0)}\) and
\(z_{\mathrm{fold}}^{(1)}\) into \(\nu\) vectors, and split
\(v_{\mathrm{aux}}\) into \(\mu_{\mathrm{split}}\) vectors.  After zero
padding, the next relation has
\begin{equation}
 n'=\max\left\{\left\lceil\frac{n}{\nu}\right\rceil,
                 \left\lceil\frac{m_{\mathrm{aux}}}{\mu_{\mathrm{split}}}
                 \right\rceil\right\},
 \qquad
 r'=2\nu+\mu_{\mathrm{split}}.
 \label{eq:next-rank-multiplicity}
\end{equation}
Parameters are chosen so that
\(n/\nu\approx m_{\mathrm{aux}}/\mu_{\mathrm{split}}\), avoiding excessive
zero padding.

Let the coefficient standard deviation of the current witness be
\(s=\beta/\sqrt{rnd}\), and let
\(\tau_{\LaB}=71\).  The heuristic parameter rules of \cite{BS23} are
\begin{align}
 b&\approx\left\lfloor\sqrt{\sqrt{12r\tau_{\LaB}}\,s}\right\rceil,
 \label{eq:base-b}\\
 t_1&=\left\lceil\frac{\log\proofmod}{\log b}\right\rceil,
 \qquad b_1=\left\lceil\proofmod^{1/t_1}\right\rceil,
 \label{eq:t1}\\
 t_2&=\left\lceil
 \frac{\log(\sqrt{24nd\,s^2})}{\log b}
 \right\rceil,
 \qquad
 b_2=\left\lceil(\sqrt{24nd\,s^2})^{1/t_2}\right\rceil,
 \label{eq:t2}\\
 \gamma&=\beta\sqrt{\tau_{\LaB}},
 \label{eq:gamma-lab}\\
 \gamma_1^2&=\frac{b_1^2t_1}{12}r\kappa d
       +\frac{b_2^2t_2}{12}\frac{r^2+r}{2}d,
 \label{eq:gamma1-lab}\\
 \gamma_2^2&=\frac{b_1^2t_1}{12}\frac{r^2+r}{2}d,
 \label{eq:gamma2-lab}\\
 (\beta')^2&=\frac{2}{b^2}\gamma^2+\gamma_1^2+\gamma_2^2.
 \label{eq:beta-prime}
\end{align}
The heuristic quantities must be checked against the vectors actually produced
by the implementation.  If a vector exceeds the predicted norm, the prover
must restart or increase the commitment rank/bounds; soundness must never rely
only on the heuristic estimate.

The final recursion level omits outer commitments and uses the interactive
``fewer garbage polynomials'' optimization, because all openings are sent at
that level.  This is the reason the final level contributes most of the proof
bytes in the concrete schedule.

% =============================================================================
\section{Composition, exactness, and predicate special soundness}
% =============================================================================

Define the interactive composed protocol \(\LNPLab\) as follows:
\begin{enumerate}
  \item run \(\LNPcom\) through the point at which its final witness
  \(w_{\mathrm{bdry}}\) would be sent;
  \item instead of sending \(w_{\mathrm{bdry}}\), run recursively composed
  LaBRADOR to prove
  \((x_{\mathrm{bdry}},w_{\mathrm{bdry}})\in\calR_{\mathrm{bdry}}\);
  \item apply Fiat--Shamir to the complete public-coin composition.
\end{enumerate}
The final NIZK is
\[
  \FS[\LNPLab].
\]

Both LNP and one LaBRADOR level are nine-message public-coin protocols.  In the
coordinate-wise predicate-special-soundness notation used by the formal
analysis, a LaBRADOR level has
\[
  K=(2,2,2,3),\qquad R=(1,1,1,r).
\]
After \(t\) recursive levels, the predicate system is the LNP predicate system
followed by \(t\) copies of the LaBRADOR predicate system.  The binding branch
is a union of Module-SIS relations arising from the LNP commitments, the three
compressing commitments, and each LaBRADOR level.

The intended application witness remains \emph{exact}: when extraction follows
the non-binding branch, LNP recovers a witness satisfying the original norm
bounds and integer no-wrap conditions.  LaBRADOR's JL norm slack affects the
short vectors appearing in the alternative binding/MSIS relations and the
recursive proof witness; it does not replace the exact application relation by
a relaxed one.

For \(N\) accepting proofs, the multi-proof extractor loses linearly in \(N\):
if \(\kappa\) is the single-proof knowledge error of the full predicate system,
then
\[
  \kappa_{\LNPLab,N}=N\kappa.
\]

% =============================================================================
\section{Concrete recursive schedule and draft proof-size estimate}
% =============================================================================

The current thesis reports the following schedule for a boundary circuit below
\(2^{23}\) binary constraints.

\begin{center}
\begin{tabular}{@{}rrrrr@{}}
\toprule
Level & Rank \(n\) & Multiplicity \(r\) & Witness size (KiB) & Output (KiB)\\
\midrule
R1CS \(\to\calR_{\LaB}\) & -- & -- & 1024.00 & 0.6250\\
1 & 27,595 & 38 & 8192.00 & 4.2046\\
2 & 3,450 & 21 & 2235.37 & 4.2037\\
3 & 1,150 & 10 & 423.65 & 3.5613\\
4 & 575 & 7 & 147.75 & 3.5372\\
5 & 290 & 7 & 73.27 & 3.5191\\
6 & 290 & 5 & 53.13 & 4.1388\\
7 & 153 & 4 & 40.39 & 30.7189\\
\midrule
\multicolumn{4}{r}{Recursive contribution} & 54.5087\\
\bottomrule
\end{tabular}
\end{center}

With \(\kappa_{\mathrm{MSIS}}=2\), the displayed auxiliary communication is
\begin{align*}
  3\text{ compact transcript commitments}
    &=3\cdot640\text{ bytes}=1.875\text{ KiB},\\
  4\text{ constant-term mask commitments}
    &=4\cdot320\text{ bytes}=1.250\text{ KiB},\\
  z_{\mathrm{JL}}
    &=608\text{ bytes}=0.59375\text{ KiB}.
\end{align*}
Therefore the draft estimate is
\begin{equation}
  \abs{\pi_{\LNPLab}}
  \approx54.5087+1.875+1.250+0.59375
  =58.22745\text{ KiB}.
  \label{eq:draft-proof-size}
\end{equation}
This accounting treats all public-coin challenges as 128-bit seeds expanded by
Fiat--Shamir and follows the thesis's current boundary-circuit estimate of
\(7{,}600{,}000<2^{23}\) binary constraints.

\begin{importantbox}[Status of the \(58.23\) KiB number]
The number in \cref{eq:draft-proof-size} is a useful engineering target, not a
final theorem-level concrete size.  It must be regenerated after:
(i) serializing both halves \(z_L,z_R\),
(ii) fixing all four norm slacks,
(iii) fixing the complete quotient/carry bounds,
(iv) fixing the real projected bit-block dimension, and
(v) selecting \(\kappa_A\), the LNP challenge parameters, and every
LaBRADOR per-level Module-SIS rank with a concrete estimator.
\end{importantbox}

% =============================================================================
\section{Implementation checklist}
% =============================================================================

Before treating the transcript as fully instantiated, the implementation should
complete the following checks.

\begin{enumerate}
  \item \textbf{Canonical relation.} Confirm that the implementation proves
  \cref{eq:R1-native,eq:R2-native,eq:R3-native,eq:R4-native} with the full
  \(z=(z_L,z_R)\).
  \item \textbf{Dimensions.} Generate every message length from a single typed
  relation specification; do not copy the inconsistent summary-table entries.
  \item \textbf{Quotients.} Derive coefficient bounds for all \(1{,}458\)
  quotient/carry outputs, encode them, and include their reconstruction and
  bitness constraints.
  \item \textbf{Projection serialization.} List every ring coordinate in
  \(x_{\mathrm{short}}\), compute \(D_{\mathrm{proj}}\) and
  \(\mu_{\mathrm{proj}}\), and rerun \cref{eq:sproj,eq:no-wrap-summary}.
  \item \textbf{JL distribution.} Use centered \(\Bin_1\) entries, or reprove
  the constants for another distribution.
  \item \textbf{LNP challenge.} Fix
  \((\kappa_{\mathrm{ch}},k_{\mathrm{ch}},\eta_{\mathrm{ch}})\), measure its
  acceptance probability, prove challenge-difference invertibility, and update
  Gaussian and MSIS bounds.
  \item \textbf{Commitment security.} Select \(\kappa_A\), opening widths, and
  all compact-commitment ranks using Module-LWE/Module-SIS estimators.
  \item \textbf{Boundary completeness.} Ensure the LaBRADOR circuit checks
  projection, constant coefficients, challenge recomputation, parsing,
  \(f_1,f_2,f_3\), response norm, and all three compact openings.
  \item \textbf{Post-aggregation factorization.} Build the optimized quadratic
  circuit only after \(\boldsymbol\mu\) is known; compare it symbolically and
  with randomized tests to the unfactored \(Q_{\boldsymbol\mu}\).
  \item \textbf{Recursion.} Generate \((n_i,r_i,b_i,b_{1,i},b_{2,i},t_{1,i},
  t_{2,i},\kappa_i,\kappa_{1,i},\kappa_{2,i})\) from code, then recompute every
  proof-size row rather than hard-coding the table.
  \item \textbf{Fiat--Shamir.} Domain-separate every challenge and bind the CRS,
  complete statement, preceding messages, dimensions, and protocol version.
  \item \textbf{Test vectors.} Publish at least one accepted transcript and
  negative tests that alter each committed flow, projection entry, constant
  coefficient, response, and dimension field.
\end{enumerate}

% =============================================================================
\section{Compact protocol summary}
% =============================================================================

The complete proof pipeline is
\[
 \boxed{
 \begin{array}{c}
 \text{native NIBS relation over }R_{q_{\mathrm{sig}}}
 \\ \Downarrow\ \text{bounded quotient/carry lift}
 \\ \text{exact relation over }R_{q_\pi}
 \\ \Downarrow\ \LNPcom\ \text{with three compressed transcript flows}
 \\ \calR_{\mathrm{bdry}}
 \\ \Downarrow\ \text{binary-R1CS / dot-product reduction}
 \\ \LaB^{\circ t}
 \\ \Downarrow\ \text{Fiat--Shamir}
 \\ \FS[\LNPLab]
 \end{array}}
\]
LNP supplies zero knowledge and exact extraction for the application witness;
LaBRADOR supplies recursive succinctness for the large final opening relation.
The three compact transcript commitments are necessary because compressing only
the final LNP responses would leave statement-dependent intermediate flows in
the public proof.

\begin{thebibliography}{9}

\bibitem{LNP22}
V. Lyubashevsky, N. K. Nguyen, and M. Plancon.
\newblock \emph{Lattice-Based Zero-Knowledge Proofs and Applications: Shorter,
Simpler, and More General}.
\newblock IACR Cryptology ePrint Archive, Report 2022/284, 2022.

\bibitem{BS23}
W. Beullens and G. Seiler.
\newblock \emph{LaBRADOR: Compact Proofs for R1CS from Module-SIS}.
\newblock IACR Cryptology ePrint Archive, Report 2022/1341; CRYPTO 2023.

\bibitem{DPKNRT25}
R. del Pino, S. Katsumata, G. Niot, M. Reichle, and K. Takemure.
\newblock \emph{Unmasking TRaccoon: A Lattice-Based Threshold Signature with an
Efficient Identifiable Abort Protocol}.
\newblock IACR Cryptology ePrint Archive, Report 2025/849, 2025.

\bibitem{AAB24}
M. A. Aardal et al.
\newblock \emph{Aggregating Falcon Signatures with LaBRADOR}.
\newblock CRYPTO 2024.

\end{thebibliography}

\end{document}
